\section{Scientific scope}
\label{sec:introduction}

PULSUS calculates linear and third-order response functions for open
quantum systems. It evaluates frequency-domain spectra directly from the
field-free Liouvillian, so a time-domain interferogram is not needed first.

\subsection{Supported systems}
Provide the Hamiltonian, dipole and collapse operators as finite matrices.
The package supports coupled electronic models and other systems that can
be represented in a finite basis. For a Hilbert-space dimension $d$, each
operator is $d\times d$ and the dense Liouvillian is $d^2\times d^2$.
The matrix size therefore sets the computational cost; selected-coordinate
calculations reduce the number of spectral points evaluated.

\subsection{Assumptions of the finite-pulse calculation}
The standard finite-pulse formulation uses a time-independent Markovian
Lindblad generator, electric-dipole coupling, Gaussian pulses with ordered,
separated interactions, and a stationary initial state under the field-free
Liouvillian. Use \texttt{system.stationary\_state()} to obtain that state and
pass it explicitly as \texttt{rho0}. Although several functions accept
other density matrices, accepting a matrix does not establish the validity
of the standard finite-pulse formula for a nonstationary preparation.

The public time-domain functions implement impulsive interactions.
Finite-duration driven time propagation, non-Markovian dynamics,
time-dependent field-free generators, and automatic construction of
molecular Hamiltonians or dipole components are outside the documented
production interface.

\subsection{Calculation workflow}
\begin{enumerate}
\item Construct the system in a consistent unit system and prepare its state.
\item Define Gaussian pulses for a finite-pulse or short-pulse model.
\item Select a response model, R/NR convention, waiting time and frequency axes.
\item Calculate the complex response and display the desired component.
\end{enumerate}
